\emph{Preliminaries.} Only clauses (b), (c) and the torsion equation of the Definition are used; its field equation is not. From the Definition: $D_{a}M = a \cdot \nabla M + \Omega(a) \times M$ and $A \cdot D = D_{\hb(A)}$ (clause (b)); $[D_{a}, D_{b}]M = R(a \wedge b) \times M$ for constant $a, b$, with $\mathcal{R}(B) = -R(\hb(B))$ (clause (c)); and the torsion equation $D \wedge \hb^{T}(a) = \hb^{T}(\nabla \wedge a)$ (clause (f)). $\Omega(a) \times$ is a derivation of the product and $(W \times X) \cdot Y + X \cdot (W \times Y) = 0$ for a bivector $W$ and vectors $X, Y$, so $A \cdot D$ is a derivation of every product of multivector fields and $\hb(A) \cdot \nabla(X \cdot Y) = (A \cdot D\,X) \cdot Y + X \cdot (A \cdot D\,Y)$. For a bivector $\mathcal{B}$ and vectors $x, y$: $x \cdot (\mathcal{B} \cdot y) = -y \cdot (\mathcal{B} \cdot x)$. For physical vector fields put $T(A, C) = A \cdot D\,C - C \cdot D\,A$.

\emph{Lemma 1 (the bracket).} $\hb\bigl(T(A, C)\bigr) = (a \cdot \nabla)c - (c \cdot \nabla)a$ for $a = \hb(A)$, $c = \hb(C)$: the bracket of the background fields. Proof: for physical vectors $A, C$ and a vector field $X$, $(A \wedge C) \cdot (D \wedge X) = A \cdot (C \cdot D\,X) - C \cdot (A \cdot D\,X)$, since $D \wedge X = \partial_{E} \wedge (E \cdot D\,X)$ and $(A \wedge C) \cdot (u \wedge v) = (A \cdot v)(C \cdot u) - (A \cdot u)(C \cdot v)$. The torsion equation at the constant vector $e^{\mu}$ gives $D \wedge g^{\mu} = 0$ for $g^{\mu} = \hb^{T}(e^{\mu})$, so $A \cdot (C \cdot D\,g^{\mu}) = C \cdot (A \cdot D\,g^{\mu})$. The scalar $C \cdot g^{\mu} = c \cdot e^{\mu}$ has derivative $(a \cdot \nabla)(c \cdot e^{\mu}) = (A \cdot D\,C) \cdot g^{\mu} + C \cdot (A \cdot D\,g^{\mu})$; subtracting the same with $A$ and $C$ exchanged gives $\bigl((a \cdot \nabla)c - (c \cdot \nabla)a\bigr) \cdot e^{\mu} = T(A, C) \cdot g^{\mu} = \hb(T(A, C)) \cdot e^{\mu}$ for each $\mu$. Consequently $T(A, T(B, C)) + T(B, T(C, A)) + T(C, T(A, B)) = 0$: its image under $\hb$ is the cyclic sum of iterated brackets of the background fields $a, b, c$, which vanishes because second derivatives commute.

\emph{Lemma 2 (the commutator).} $[A \cdot D, C \cdot D]X = -\mathcal{R}(A \wedge C) \times X + T(A, C) \cdot D\,X$ for a multivector field $X$. Proof: $A \cdot D(C \cdot D\,X) = D_{a}(D_{c}X)$ with $c$ a field; $D_{c}$ is linear in $c$, so this is $D_{a}D_{c}X$ with $c$ held fixed plus $D_{(a \cdot \nabla)c}X$. Hence $[A \cdot D, C \cdot D]X = R(a \wedge c) \times X + D_{(a \cdot \nabla)c - (c \cdot \nabla)a}X$; by Lemma 1 the last term is the derivative along $T(A, C)$, and $R(a \wedge c) = R(\hb(A \wedge C)) = -\mathcal{R}(A \wedge C)$.

\emph{(i), first identity.} By Lemma 2 on a vector field, $[A \cdot D, B \cdot D]C - T(A, B) \cdot D\,C = -\mathcal{R}(A \wedge B) \cdot C$. Sum over the cyclic permutations of $A, B, C$. The left side is $\sum\bigl(A \cdot D(B \cdot D\,C) - B \cdot D(A \cdot D\,C) - T(A, B) \cdot D\,C\bigr)$; relabelling the second and third terms cyclically it is $\sum\bigl(A \cdot D(B \cdot D\,C - C \cdot D\,B) - T(B, C) \cdot D\,A\bigr) = \sum T(A, T(B, C))$, which is zero by Lemma 1. So the cyclic sum of $\mathcal{R}(A \wedge B) \cdot C$ vanishes for vector fields, and both sides being linear over the scalar functions, for vectors at a point.

\emph{(i), second identity.} For three operators the cyclic sum of $[A \cdot D, [B \cdot D, C \cdot D]]$ vanishes. By Lemma 2 and the derivation property, on a vector field $X$,
\[
[A \cdot D, [B \cdot D, C \cdot D]]X = -\bigl(A \cdot D(\mathcal{R}(B \wedge C))\bigr) \times X - \mathcal{R}(A \wedge T(B, C)) \times X + T(A, T(B, C)) \cdot D\,X ,
\]
the first term from differentiating $-\mathcal{R}(B \wedge C) \times X$ and the other two from $[A \cdot D, T(B, C) \cdot D]X$. The cyclic sum of the last term vanishes by Lemma 1. So $W \times X = 0$ for every vector field $X$, with $W$ the cyclic sum of $A \cdot D(\mathcal{R}(B \wedge C)) + \mathcal{R}(A \wedge T(B, C))$; a bivector whose inner product with every vector vanishes is zero, so $W = 0$. Now $A \cdot D(\mathcal{R}(B \wedge C)) = (A \cdot D\,\mathcal{R})(B \wedge C) + \mathcal{R}((A \cdot D\,B) \wedge C) + \mathcal{R}(B \wedge (A \cdot D\,C))$ and $A \wedge T(B, C) = A \wedge (B \cdot D\,C) - A \wedge (C \cdot D\,B)$. In the cyclic sum the twelve bivectors under $\mathcal{R}$ cancel in pairs: $(A \cdot D\,B) \wedge C$ with $C \wedge (A \cdot D\,B)$, $(B \cdot D\,C) \wedge A$ with $A \wedge (B \cdot D\,C)$, $(C \cdot D\,A) \wedge B$ with $B \wedge (C \cdot D\,A)$, and $B \wedge (A \cdot D\,C)$, $C \wedge (B \cdot D\,A)$, $A \wedge (C \cdot D\,B)$ each with its negative. What remains is the second identity.

\emph{(ii).} A linear map of the bivectors has $36$ coefficients on the unit bivectors $e_{i}e_{j}$; the first identity on the four triples of frame vectors is $16$ linear conditions, of rank $16$ (computed), so $\mathcal{K}$ has dimension $20$. Each element of a basis of $\mathcal{K}$ is symmetric for the product of bivectors and has a symmetric Ricci map (computed on the basis; both properties are linear). The identity is formed with the product, so the rotors carry $\mathcal{K}$ into itself (checked for the six unit bivectors on a sample of the basis). For the last sentence: the field $\hb^{-1} = \id + \tfrac14\sum_{j}(p_{j} \cdot x)^{2}h_{j}$ has at $x = 0$, where $\hb$ is the identity and $\Omega$ vanishes, the curvature operator $\sum_{j}\mathcal{R}_{p_{j}}[h_{j}]$, the stored first-order curvatures, by the last statement of clause (iv) of Theorem \claimref{synthesis.null-plane-data-of-the-first-order-gravitational-field-are-the-shear-of-the-generators} (its proof is linear in the second derivatives of the field at that point); these lie in $\mathcal{K}$ by (i), and for ten momenta and the ten unit symmetric maps they span a space of dimension $20$ (computed), so they span $\mathcal{K}$.

\emph{(iii).} The derivative $A \cdot D$ respects the products, so it passes through the contraction with a frame: the Ricci map of $A \cdot D\,\mathcal{R}$ is $A \cdot D$ of the Ricci map, and the scalar of that is $\hb(A) \cdot \nabla$ of the scalar. In the second identity take $A = \gamma_{a}$, $B = \gamma_{b}$ of a frame and apply $\gamma^{b} \cdot (\gamma^{a} \cdot\,\cdot\,)$, summed. The third term gives $\gamma^{b} \cdot \bigl(\gamma^{a} \cdot (C \cdot D\,\mathcal{R})(\gamma_{a} \wedge \gamma_{b})\bigr)$, the scalar of $C \cdot D\,\mathcal{R}$, that is $C \cdot D\,\mathcal{R}$ with $\mathcal{R}$ the scalar. The second gives $-\gamma^{b} \cdot \bigl(\gamma^{a} \cdot (\gamma_{b} \cdot D\,\mathcal{R})(\gamma_{a} \wedge C)\bigr) = -\gamma^{b} \cdot \bigl((\gamma_{b} \cdot D\,\mathcal{R})(C)\bigr)$, minus the divergence of the Ricci map at $C$. The first gives $\gamma^{b} \cdot \bigl(\gamma^{a} \cdot (\gamma_{a} \cdot D\,\mathcal{R})(\gamma_{b} \wedge C)\bigr) = -\gamma^{a} \cdot \bigl(\gamma^{b} \cdot (\gamma_{a} \cdot D\,\mathcal{R})(\gamma_{b} \wedge C)\bigr)$, since $x \cdot (y \cdot \mathcal{B}) = -y \cdot (x \cdot \mathcal{B})$ for a bivector, again minus the divergence of the Ricci map at $C$. So $2\operatorname{div}\mathcal{R}(C) = C \cdot D\,\mathcal{R}$. For the identity times a scalar field $f$, $(A \cdot D\,(f\,\id))(C) = (\hb(A) \cdot \nabla f)\,C$, so its divergence at $C$ is $\hb(C) \cdot \nabla f = C \cdot D f$. Hence $\operatorname{div}\mathcal{G}(C) = \tfrac12 C \cdot D\,\mathcal{R} - \tfrac12 C \cdot D\,\mathcal{R} = 0$. The check verifies both identities and the contraction exactly on a field that is neither weak nor symmetric, with $\Omega$ solved from the torsion equation: the cyclic identity on all triples of the frame, the differential identity on the same triples with all twelve of its terms nonzero, $\operatorname{div}\mathcal{R} = \tfrac12 D\mathcal{R}$ in the four directions with the right side nonzero in each, and $\operatorname{div}\mathcal{G} = 0$.

\emph{(iv).} A map $\Phi$ commutes with the rotors exactly when it commutes with their generators, the bivectors, acting by the commutator product: on a curvature operator by $(G \cdot \mathcal{R})(B) = G \times \mathcal{R}(B) - \mathcal{R}(G \times B)$ and on a map of vectors by $(G \cdot F)(a) = G \times F(a) - F(G \times a)$. In a basis of $\mathcal{K}$, $\Phi$ has $20 \times 16 = 320$ coefficients and the six unit bivectors give a linear system in them with rational coefficients. Its rank modulo the prime $1000003$ is $318$ (computed; the denominators are prime to it); a minor that is nonzero modulo a prime is nonzero, so the rational rank is at least $318$ and the solutions form a space of dimension at most two. The Ricci map and the identity times the scalar are formed with the product, so they commute with the rotors (verified: they satisfy every equation of the system exactly), and they are independent on $\mathcal{K}$ (computed). So the space has dimension two. The same system for the maps of vectors alone, $G \cdot F = 0$, has a one-dimensional space of solutions, the multiples of the identity (computed). An affine function of the curvature operator is $F_{0} + \Phi(\mathcal{R})$ with $\Phi$ linear; by (ii) the curvature operator at a point ranges over all of $\mathcal{K}$, zero included, so covariance under the rotation gauge at $\mathcal{R} = 0$ makes $F_{0}$ commute with the rotors, and then $\Phi$ does: $F_{0} = \lambda\,\id$ and $\Phi(\mathcal{R})(a) = \alpha\,\mathcal{R}(a) + \beta\,a\,\mathcal{R}$. That $\alpha, \beta, \lambda$ are constants and not functions of position, of $\hb$ or of $\Omega$ is the hypothesis that $F$ is one and the same function of the curvature operator at every point and of nothing else; without it the divergence in (v) would acquire the gradients of the coefficients.

\emph{(v).} By (iii), $\operatorname{div}(\alpha\,\mathcal{R}(\cdot))(C) = \tfrac12\alpha\,C \cdot D\,\mathcal{R}$, $\operatorname{div}(\beta\,\mathcal{R}\,\id)(C) = \beta\,C \cdot D\,\mathcal{R}$ and $\operatorname{div}(\lambda\,\id) = 0$. The trace of the Ricci map is the scalar and that of the identity is four. If $\tfrac12\alpha + \beta \neq 0$ the divergence is nonzero wherever the scalar is not constant, and such fields exist (the field of the check has $D\mathcal{R} \neq 0$ in all four directions).

\emph{(vi).} Take the divergence of the equation: by (v) and (C), $(\tfrac12\alpha + \beta)\,C \cdot D\,\mathcal{R} = 0$ for every $C$ on every solution. Take its trace: $(\alpha + 4\beta)\,\mathcal{R} + 4\lambda = -\kappa\,\partial_{a} \cdot T(a)$. Suppose $\tfrac12\alpha + \beta \neq 0$. Then the gradient of the scalar $\mathcal{R}$ vanishes on every solution, and so does the gradient of the trace of the stress: if $\alpha + 4\beta \neq 0$ because it is $-(\alpha + 4\beta)\mathcal{R}/\kappa - 4\lambda/\kappa$, and if $\alpha + 4\beta = 0$ because it is $-4\lambda/\kappa$. That contradicts (C), which asks for matter whose trace has a gradient that does not vanish identically. So $\beta = -\tfrac12\alpha$ and $F = \alpha\,\mathcal{G} + \lambda\,\id$. If $\alpha = 0$ the equation is $\lambda\,a = -\kappa\,T(a)$, which no matter with a trace of nonvanishing gradient satisfies. Dividing by $\alpha$ gives the stated equation.

\emph{(vii).} For a weak field $\hb^{-1} = \id + s$ with $s$ symmetric, the torsion equation gives $\Omega(a) = -\nabla \wedge s(a)$ to first order, the derivative acting on $s$, and the curvature operator is $\mathcal{R}(a \wedge b) = -\bigl((a \cdot \nabla)\Omega(b) - (b \cdot \nabla)\Omega(a)\bigr) = -\tfrac12\,\nabla \wedge h((a \wedge b) \cdot \nabla)$ with $h = 2s$ the first-order part of the derived metric: the stored first-order curvature, as in the proof of clause (iv) of the cited theorem. So $\mathcal{G}$ is at first order the stored $\mathcal{E}[h]$, and physical and background vectors agree at this order. The check confirms it on the field $\hb^{-1} = \id + \tfrac14(p \cdot x)^{2}h_{0}$ at $x = 0$, comparing the Einstein map computed from the gauge fields with the stored first-order Einstein map on the four frame vectors. With $\Lambda = 0$ the equation of (vi) is then $\mathcal{E}[h] = -(\kappa/\alpha)\,T$ at first order. Its operator is the Einstein map, which is homogeneous of second order and covariant, vanishes on displacements and is symmetric for the pairing, the conditions (L), (D), (V) of clause (v) of the cited theorem; by its clause (vi), with the response of matter of Theorem \claimref{synthesis.weyl-flat-variations-carry-only-the-stress-trace-and-the-whole-stress-clause-is-the-trace-free-channel}(v), under its hypothesis (N) the interaction of two bodies at rest is the stored Newtonian coupling exactly when the equation is $\mathcal{E}[h] + 8\pi G\,T = 0$. So $\kappa/\alpha = 8\pi G$. The ratio is identified at $\lambda = 0$, where the flat field is the solution without matter, which is the case the stored Newtonian coupling speaks of. $\kappa/\alpha$ and $\Lambda$ are independent constants of the family of (vi), so that the same value holds for $\Lambda \neq 0$ is a convention for what $G$ names there (Newton's constant where $\Lambda$ times the square of the distances is negligible), not a deduction. This clause has no check of its own beyond the first-order comparison; it rests on the proof and on clause (vi) of the cited theorem.
